\clearpage
\thispagestyle{empty}
\pdfbookmark[0]{Der magnetische Südpol. Ausgabe und Quellen}{wilkes1841-edition}
\begin{center}\Large\textbf{Zu Wilkes’ Beobachtungen in der Nähe des magnetischen Südpols}\end{center}

Der vollständige Auszug aus Carl Friedrich Gauss’ Brief an Heinrich Christian Schumacher in den \textit{Werken}, Band V (Göttingen, 1867), S.~580, liegt hier als KI-Übertragung in modernes Deutsch mit KI-erstelltem Satz vor. Er enthält Gauss’ vorläufige Bestimmung der Lage des Pols anhand der mitgeteilten Messungen von Wilkes.

Gauss datierte den Brief auf den 9. Februar 1841 in Göttingen. Der vollständige Brief erschien in den von Heinrich Christian Schumacher herausgegebenen \textit{Astronomischen Nachrichten}, Band~18, Nr.~417, Sp.~143--144 (Altona, 4. März 1841). Die Übertragung folgt dem Wortlaut und der Zahlenschreibweise des Auszugs der Werkausgabe.

\medskip\noindent
\href{https://gdz.sub.uni-goettingen.de/id/PPN236006339}{Deutsche Werkausgabe (1867)}\par
\noindent\href{https://archive.org/details/astronomischena00berlgoog/page/n80/mode/2up}{Erstdruck des Briefs (1841)}
\clearpage
